
\documentclass[conference]{IEEEtran}

% ---- packages (kept minimal to avoid IEEEtran conflicts) -------------
\usepackage{cite}
\usepackage{amsmath,amssymb}
\usepackage{graphicx}
\usepackage{tikz}
\usepackage{booktabs}
\usepackage{xcolor}
\usepackage{url}
\usepackage{algorithm}
\usepackage{algpseudocode}

% ---- method name ------------------------------------------------------
\newcommand{\method}{\textsc{Relay}}

% ---- math shorthands -------------------------------------------------
\newcommand{\zc}{\mathbf{z}^{c}}      % content (gesture/activity) factor
\newcommand{\zd}{\mathbf{z}^{d}}      % domain (environment/layout) factor
\newcommand{\bs}[1]{\boldsymbol{#1}}

\begin{document}

% Cover page, typeset in LaTeX. The manuscript below is unmodified.
\onecolumn
\thispagestyle{empty}
\null
\begin{tikzpicture}[remember picture,overlay]
  \coordinate (covercenter) at (current page.north);
  \node[anchor=north,font=\bfseries\fontsize{10}{12}\selectfont]
    at ([yshift=-0.10in]covercenter) {PROPOSAL SKRIPSI ARTIKEL ILMIAH};

  \node[anchor=north,align=center,text width=6.6in,font=\bfseries\fontsize{9}{13}\selectfont]
    at ([yshift=-0.87in]covercenter)
    {Persepsi CSI Wi-Fi Lintas Domain melalui Pemisahan Representasi Mandiri dan \textit{Spatial}};
  \node[anchor=north,font=\bfseries\itshape\fontsize{9}{12}\selectfont]
    at ([yshift=-1.28in]covercenter) {Retrieval-Augmented Generation};

  \node[anchor=north,align=center,text width=6.4in,font=\bfseries\fontsize{9}{14}\selectfont]
    at ([yshift=-2.06in]covercenter)
    {Cross-Domain Wi-Fi CSI Perception via Self-Supervised Disentanglement and Spatial Retrieval-};
  \node[anchor=north,font=\bfseries\fontsize{9}{12}\selectfont]
    at ([yshift=-2.47in]covercenter) {Augmented Generation};

  \node[anchor=north,font=\fontsize{9}{12}\selectfont]
    at ([yshift=-3.31in]covercenter) {\textbf{Topik:} Intelegensia Semu};

  \node[anchor=north,font=\fontsize{9}{12}\selectfont]
    at ([yshift=-4.10in]covercenter)
    {2702387973 / CHRISTOPHER RICKY SOEHARTONO / Computer Science};
  \node[anchor=north,font=\fontsize{9}{12}\selectfont]
    at ([yshift=-4.52in]covercenter)
    {2702363651 / FRITZ GRADIYOGA / Computer Science};
  \node[anchor=north,font=\fontsize{9}{12}\selectfont]
    at ([yshift=-4.94in]covercenter)
    {2702363765 / JESSICA NATHANIA WENARDI / Computer Science};

  % BINUS emblem and lettering — included from logo_binus.png.
  \node[anchor=north] at ([yshift=-5.72in]covercenter)
    {\includegraphics[width=6.0cm]{logo_binus.png}};

  \node[anchor=north,font=\bfseries\fontsize{9}{12}\selectfont]
    at ([yshift=-8.44in]covercenter) {BINUS University};
  \node[anchor=north,font=\bfseries\fontsize{9}{12}\selectfont]
    at ([yshift=-8.87in]covercenter) {2026};
  \node[anchor=north west,font=\fontsize{9}{12}\selectfont]
    at ([xshift=.98in,yshift=-9.51in]current page.north west) {DIPERIKSA OLEH};
  \node[anchor=north west,align=left,font=\fontsize{8}{9}\selectfont]
    at ([xshift=.73in,yshift=-10.61in]current page.north west)
    {D3557 - Ir. Andry Chowanda, S.Kom, MM,\\
     Ph.D., MBCS, CCP, CME, IPM, SMIEEE};
\end{tikzpicture}
\clearpage
\twocolumn

\title{Cross-Domain Wi-Fi CSI Perception via Self-Supervised\\
Disentanglement and Spatial Retrieval-Augmented Generation}

\author{
\setlength{\tabcolsep}{0pt}% kontrol jarak kolom manual
\begin{tabular}{@{}c@{\hspace{0.045\textwidth}}c@{\hspace{0.045\textwidth}}c@{}}

% ================= AUTHOR 1 =================
\begin{minipage}[t]{0.28\textwidth}
\centering
\vspace{0pt}
{\fontsize{9}{12}\selectfont
Christopher Ricky Soehartono\\\\[2pt]
\textit{Computer Science Department}\\\\
\textit{School of Computer Science}\\\\
\textit{Bina Nusantara University}\\\\
Jakarta, Indonesia 11480\\\\
christopher.soehartono@binus.ac.id}
\end{minipage}

&

% ================= AUTHOR 2 =================
\begin{minipage}[t]{0.28\textwidth}
\centering
\vspace{0pt}
{\fontsize{9}{12}\selectfont
Fritz Gradiyoga\\\\[2pt]
\textit{Computer Science Department}\\\\
\textit{School of Computer Science}\\\\
\textit{Bina Nusantara University}\\\\
Jakarta, Indonesia 11480\\\\
fritz.gradiyoga@binus.ac.id}
\end{minipage}

&

% ================= AUTHOR 3 =================
\begin{minipage}[t]{0.28\textwidth}
\centering
\vspace{0pt}
{\fontsize{9}{12}\selectfont
Jessica Nathania Wenardi\\\\[2pt]
\textit{Computer Science Department}\\\\
\textit{School of Computer Science}\\\\
\textit{Bina Nusantara University}\\\\
Jakarta, Indonesia 11480\\\\
jessica.wenardi@binus.ac.id}
\end{minipage}

\end{tabular}

\\[0.8cm]

% ================= AUTHOR 4 =================
\begin{minipage}[t]{\textwidth}
\centering
\vspace{0pt}
{\fontsize{9}{12}\selectfont
Ir. Andry Chowanda, S.Kom, MM, Ph.D.,\\\\[-1pt]
{MBCS, CCP, CME, IPM, SMIEEE}\\\\[2pt]
\textit{Computer Science Department}\\\\
\textit{School of Computer Science}\\\\
\textit{Bina Nusantara University}\\\\
Jakarta, Indonesia 11480\\\\
achowanda@binus.edu}
\end{minipage}
}

\maketitle

% =====================================================================
\begin{abstract}
We introduce \method{}, a cross-domain Wi-Fi channel state information (CSI) sensing
framework that couples \emph{self-supervised disentanglement} of gesture content from
environmental conditions with a \emph{spatial retrieval-augmented} memory that supplies
the model with an explicit, updatable description of the space it currently occupies.
Existing CSI sensing systems are trained and deployed in the same room: because the
received signal is a function of both the human-induced dynamics and the static
propagation environment, accuracy collapses when the layout, transmitter placement, or
furniture configuration changes, and current remedies either require labelled target data,
a retraining round, or multiple transceivers.
Our key idea is to treat the environment not as nuisance to be erased but as
\emph{retrievable context}: a self-supervised objective separates a content code from a
domain code, and at inference a geometry-tagged memory indexed by coarse spatial
descriptors conditions the classifier on the active layout, so a new room becomes a
retrieval problem rather than a retraining problem.
We evaluate under cross-environment, cross-location, and cross-orientation protocols
against domain-adaptation and domain-generalisation baselines.
\end{abstract}

\begin{IEEEkeywords}
Wi-Fi sensing, channel state information, cross-domain generalisation, self-supervised
learning, disentangled representation, retrieval-augmented generation.
\end{IEEEkeywords}

% =====================================================================
\section{Introduction}
\label{sec:intro}

Wi-Fi sensing uses channel state information (CSI) from commodity Wi-Fi devices to
recognise human activities without cameras or additional sensing hardware
\cite{Chen2022Survey,Radwan2025Survey}. Its main challenge is poor generalisation across
environments: CSI captures both human motion and the surrounding propagation conditions,
so the same activity can produce different signals when the room, device placement, or
orientation changes \cite{Chen2022Survey,Wang2025GeneralizabilitySurvey,
Zhang2026BeyondLabels}.

Existing solutions mainly use \emph{domain adaptation} (DA), \emph{domain
generalisation} (DG), few-shot learning, or self-supervised learning. DA aligns source and
target distributions but requires target-domain data
\cite{Chen2022Fidora,Prasad2024DANN}. DG and
self-supervised approaches instead learn representations that transfer across domains
\cite{Liu2024WiSR,Liu2024WiSGP,Wang2022AirFi,Xu2023ARCFi,Yang2022AutoFi,
Barahimi2024CPC,Lyons2024WiFiAct}, while few-shot methods adapt using a small number of
target examples \cite{Yin2022FewSense,Zhao2024CrossFi,Yang2023FewCS}. However, making a
representation fully invariant to the environment may discard useful contextual
information. Disentanglement provides another option by separating activity information
from environment-specific factors \cite{Kang2021ContextAware}.

Retrieval-augmented generation (RAG) offers a way to preserve and reuse such contextual
information. Wang et al.\ \cite{Wang2025GeneralizabilitySurvey} identify external
retrieval as a possible direction for Wi-Fi sensing, while RAG has already been explored
in language models \cite{lewis2020rag}, wireless systems
\cite{Nazar2024Enwar,Nazar2026Enwar2,wi_chat}, and spatial reasoning
\cite{spatial_rag}. However, existing work does not use retrieval of spatial or layout
information as the adaptation mechanism for cross-domain CSI sensing.

We therefore propose \method{}, a framework that combines self-supervised
disentanglement with spatial retrieval. \method{} separates CSI into a \emph{content}
code representing the activity and a \emph{domain} code representing environmental
conditions. At inference time, the domain code retrieves relevant entries from a
geometry-tagged memory, and the retrieved context is used to condition the classifier.
A new environment can therefore be added through memory insertion without labelled
target data or gradient-based retraining.

% =====================================================================

\section{Literature Review}
\label{sec:related}

CSI is a sampled complex channel response; its magnitude and phase respond both to human
motion and to the static geometry of the deployment. Chen et al.\ \cite{Chen2022Survey}
formalise the domain shift this creates and group remedies into five families ---
domain-invariant feature extraction, virtual sample generation, transfer learning,
few-shot learning, and big-data solutions --- and find each has distinct limitations.
Wang et al.\ \cite{Wang2025GeneralizabilitySurvey} review over 200 papers by pipeline
stage, covering domain adaptation, meta-learning, metric learning, augmentation,
cross-modal alignment, federated and continual learning, and release a dataset platform;
they list pretraining at scale, multimodal foundation models, and continual deployment as
the open directions. Radwan et al.\ \cite{Radwan2025Survey} survey the self-supervised
variants, contrasting contrastive and non-contrastive objectives with a comparative
accuracy analysis. Zhang et al.\ \cite{Zhang2026BeyondLabels} argue that the conventional
factorisation of domains (room / user / orientation) is itself the limitation, and
redefine domains to make generalisation claims measurable. We adopt that critique: our
formulation keeps the domain as an explicit, retrievable variable rather than treating all
shift as one undifferentiated nuisance.

Adaptation-based methods assume some target access. Fidora \cite{Chen2022Fidora} combines a
variational-autoencoder augmenter with a joint classification--reconstruction classifier
and reports gains for new users (+17.8\% average F1) and for a changed environment
(beating the prior state of the art by 23.1\%). AdaWiFi \cite{Zheng2025AdaWiFi} lets
multiple IoT devices sense collaboratively and reports roughly 30\% higher accuracy under
one-shot adaptation. Prasad et al.\ \cite{Prasad2024DANN} adversarially learn
domain-invariant compressed features with a gradient reversal layer. Source-free
adaptation removes the need for source data: Wi-SFDAGR \cite{Yan2025WiSFDAGR} treats
adaptation as unsupervised clustering with an attraction--dispersion network and
uncertainty-weighted neighbours, and Li et al.\ \cite{Li2026Conformal} sidestep retraining
with conformal prediction over multivariate kernel density estimates, reporting accuracy
improvements from 30\% to 74\% across six kinds of domain variation.
\method{} differs in that it neither requires unlabelled target data at
adaptation time nor retrains: it conditions on a retrieved description of the space.

Generalisation-based methods need only source data. Widar3.0 \cite{Zheng2019Widar3}
estimates body-coordinate velocity profiles as a physically domain-independent
representation and reaches 82.6--92.4\% cross-domain with one-time training. WiSR
\cite{Liu2024WiSR} randomises subcarrier-domain ``styles'' adversarially, and WiSGP
\cite{Liu2024WiSGP} perturbs inputs along subdomain directions, both targeting
multi-factor generalisation without target data. AirFi \cite{Wang2022AirFi} learns the
critical CSI part for unseen environments, while ARC-Fi \cite{Xu2023ARCFi} targets the
more realistic semi-supervised setting where unlabelled CSI is plentiful but labels are
scarce, and blocks ``shortcut learning'' with an antenna-response-consistency augmentation
that treats co-located antennas as natural views. Our augmentations are in the same spirit
as ARC-Fi's, and we are explicit that this is inherited rather than novel.

FewSense \cite{Yin2022FewSense} pretrains a feature extractor and fine-tunes it on a few
target samples, reaching 93.9/96.5/82.7\% five-shot accuracy on SignFi/Widar/Wiar with
collaborative multi-receiver fusion adding 29.2\%. CrossFi \cite{Zhao2024CrossFi} replaces
cosine similarity with an attention-based similarity network and a class-template
generator, reporting 91.72\% one-shot and 64.81\% zero-shot cross-domain. FewCS
\cite{Yang2023FewCS} clusters intra-domain prototypes and extracts cross-domain prototypes
in a shared space for edge--cloud online learning. These systems are the closest in spirit
to retrieval --- matching a query against stored exemplars --- but their memory holds
\emph{class prototypes}, not \emph{spatial descriptions}, so they cannot exploit the fact
that the room is known.

AutoFi \cite{Yang2022AutoFi} learns from randomly captured low-quality unlabelled CSI with
a geometric self-supervised objective and transfers across tasks. WiFiAct
\cite{Lyons2024WiFiAct} pairs an environment-robust preprocessing pipeline with a Bayesian
CNN and contrastive augmentation. Context-aware predictive coding \cite{Barahimi2024CPC}
and several recent masked autoencoder, JEPA, and adapter-based objectives further
demonstrate that strong CSI encoders can be built without labels. Jiang et al.\
\cite{Jiang2025Foundation} pretrain masked autoencoders on 1.3M samples from 14 datasets
and find log-linear gains with pretraining data --- concluding that data, not capacity, is
the bottleneck. \method{} inherits this encoder-scaling result: our contribution is what
happens \emph{after} the encoder, and is complementary to pretraining scale.

RAG conditions generation on retrieved non-parametric context \cite{lewis2020rag}, and
spatial variants constrain retrieval with geometric relations \cite{spatial_rag}. In
wireless, ENWAR \cite{Nazar2024Enwar} builds a RAG-empowered multimodal LLM framework for
wireless environment perception, extended in ENWAR~2.0 \cite{Nazar2026Enwar2} with adaptive
knowledge-base formation, chain-of-thought reasoning and beam tracking; Wi-Chat
\cite{wi_chat} shows an LLM can interpret CSI for activity recognition when prompted with
sensing principles. These systems put retrieval in service of language-level perception.
Wang et al.\ \cite{Wang2025GeneralizabilitySurvey} raise the open question of RAG as
on-demand Wi-Fi knowledge. To our knowledge, no prior work uses a spatially-indexed
retrieval memory as the cross-domain adaptation mechanism for a CSI sensing head, which is
the gap we target. 
% =====================================================================
\section{Method}
\label{sec:method}

\begin{table}[t]
\centering
\caption{Notation used in Section~\ref{sec:method}.}
\label{tab:notation}
\small
\begin{tabular}{p{0.22\columnwidth}p{0.65\columnwidth}}
\toprule
\textbf{Symbol} & \textbf{Meaning} \\
\midrule
$x$ & CSI input window. \\
$N_s, N_t$ & Number of subcarriers and time samples in a CSI window. \\
$y, C$ & Activity label and number of activity classes. \\
$d$ & Domain descriptor available at inference time. \\
$E_\theta$ & Encoder that maps a CSI window to content and domain codes. \\
$D_\psi$ & Decoder used for CSI reconstruction. \\
$g_\phi$ & Activity-classification head. \\
$\zc, \zd$ & Content code and domain/environment code, respectively. \\
$\mathcal{M}$ & Spatially indexed retrieval memory. \\
$\rho_m, e_m$ & Geometric descriptor and conditioning embedding of memory entry $m$. \\
$\phi_{\mathrm{geom}}(\cdot)$ & Function that extracts a geometric descriptor from CSI. \\
$K, S$ & Number of retrieved entries and the resulting top-$K$ set. \\
$s(\cdot,\cdot)$ & Similarity function used for memory retrieval. \\
$w_m$ & Softmax retrieval weight assigned to memory entry $m$. \\
$\tau$ & Softmax temperature for retrieval weighting. \\
$\delta$ & Similarity threshold for the out-of-memory decision. \\
$\mathcal{L}_{\mathrm{rec}}$ & Reconstruction loss. \\
$\mathcal{L}_{\mathrm{cls}}$ & Activity-classification loss. \\
$\mathcal{L}_{\mathrm{ind}}$ & Independence/disentanglement loss. \\
$\mathcal{L}_{\mathrm{con}}$ & Content-consistency loss across augmented views. \\
$\lambda_{1..3}$ & Weights of the classification, independence, and consistency losses. \\
\bottomrule
\end{tabular}
\end{table}

\subsection{Problem formulation}
Let $x\in\mathbb{R}^{N_s\times N_t}$ denote a CSI window with $N_s$ subcarriers and $N_t$
time samples, and $y\in\{1,\dots,C\}$ its activity label. We associate each measurement
with a domain descriptor $d$ that is available \emph{at inference} without labels: a coarse
geometric signature of the space (transmitter--receiver geometry, coarse distance bins
derived from the signal statistics) optionally plus installation metadata. Standard cross-domain sensing trains
$f_\theta(x)\!\to\!y$ on source domains $\{d_s\}$ and is evaluated on an unseen $d_t$,
treating $d$ as nuisance. We instead model
\begin{equation}
p(y\mid x, \mathcal{M}) \;=\; \sum_{m \in \mathrm{TopK}(x,\mathcal{M})} w_m \, p_\phi\!\left(y \mid \zc(x),\, m\right),
\label{eq:main}
\end{equation}
where $\zc(x)$ is a content code, $\mathcal{M}=\{(\rho_m, e_m)\}$ is a memory of
spatial/layout entries ($\rho_m$ a geometric descriptor, $e_m$ its conditioning embedding),
and $w_m$ are retrieval weights (Eq.~\eqref{eq:main}). Adaptation to a new environment
consists of inserting an entry into $\mathcal{M}$; the parameters $\theta,\phi$ are untouched.

\subsection{Stage 1: self-supervised disentanglement}
An encoder $E_\theta$ maps $x$ to $(\zc,\zd)$ by splitting the embedding into two
sub-vectors, decoded jointly:
\begin{equation}
\hat{x} = D_\psi(\zc, \zd), \qquad
\mathcal{L}_{\mathrm{rec}} = \|x-\hat{x}\|_2^2 .
\end{equation}
Content must be predictive of the label for labelled source windows,
$\mathcal{L}_{\mathrm{cls}} = \mathrm{CE}\big(y, g_\phi(\zc)\big)$, while the domain code
must be predictable from the signal statistics but \emph{not} from the content code, and
vice versa. We enforce statistical independence with an HSIC-style penalty:
\begin{equation}
\mathcal{L}_{\mathrm{ind}} = \mathrm{HSIC}(\zc, \zd) \;\; \text{or} \;\;
\min_{E_\theta}\max_{h}\ \mathcal{L}_{\mathrm{adv}}(h(\zc),\zd).
\end{equation}
Finally we add physics-informed augmentations that generate views of the same event which
differ in domain but not in content --- co-located antenna responses
\cite{Xu2023ARCFi}, subcarrier-domain style randomisation \cite{Liu2024WiSR}, and
subdomain-guided perturbation \cite{Liu2024WiSGP} --- with a consistency loss
$\mathcal{L}_{\mathrm{con}}$ between the content codes of two views. The total source
objective is
\begin{equation}
\mathcal{L} = \mathcal{L}_{\mathrm{rec}} + \lambda_1 \mathcal{L}_{\mathrm{cls}}
+ \lambda_2 \mathcal{L}_{\mathrm{ind}} + \lambda_3 \mathcal{L}_{\mathrm{con}} .
\label{eq:total}
\end{equation}

\subsection{Stage 2: spatially-indexed retrieval memory}
Each memory entry stores a geometric descriptor and the domain embedding of a reference
capture from that space,
$\rho_m = \big[\phi_{\mathrm{geom}}(x^{(m)}_{\mathrm{ref}}),\, \text{meta}_m\big]$,
$e_m = \zd(x^{(m)}_{\mathrm{ref}})$. At inference we retrieve the top-$K$ entries by
descriptor similarity, weight them by softmax over similarity, and condition the label
head on the concatenation $[\zc(x); \sum_m w_m e_m]$. Because $e_m$ is an embedding, a new
room can also be added \emph{label-free}: a short unattended capture supplies $\rho_m$ and
$e_m$. Algorithm~\ref{alg:relay} states the procedure.
\begin{equation}
w_m = \frac{\exp\!\big(s(\rho(x),\rho_m)/\tau\big)}{\sum_{k\in\mathrm{TopK}}\exp\!\big(s(\rho(x),\rho_k)/\tau\big)} .
\end{equation}

\begin{algorithm}[t]
\caption{\method{} inference (adaptation = memory insertion, no retraining)}
\label{alg:relay}
\begin{algorithmic}[1]
\Require frozen $E_\theta$, head $g_\phi$, memory $\mathcal{M}$, window $x$ from an unseen space
\State $(\zc, \zd) \gets E_\theta(x)$
\State $\rho \gets \phi_{\mathrm{geom}}(x)$
\State $S \gets \mathrm{TopK}\big(s(\rho,\rho_m)\ \forall m\big)$
\If{$\max_m s(\rho,\rho_m) < \delta$}                       \Comment{unknown space}
    \State flag \textsc{out-of-memory}; optionally insert a label-free entry from a short capture
\EndIf
\State $w_m \gets \mathrm{softmax}\big(s(\rho,\rho_m)/\tau\big)$ over $S$
\State \Return $g_\phi\big([\zc; \textstyle\sum_{m\in S} w_m e_m]\big)$
\end{algorithmic}
\end{algorithm}

\subsection{Why retrieval instead of pure invariance}
Under cross-entropy with a domain-adversarial penalty the encoder is pushed to remove
domain information. If two domains require different decision boundaries --- which is the
empirical finding behind the multi-subdomain results of \cite{Liu2024WiSGP} and the
domain-redefinition argument of \cite{Zhang2026BeyondLabels} --- invariance removes the
information needed to select the right boundary. Keeping $\zd$ and \emph{retrieving} the
matching context makes the boundary a function of retrieved state, which is exactly what
the deployment regime offers for free: an installer who knows the room. This also gives the
system an honest failure mode via the out-of-memory branch in
Algorithm~\ref{alg:relay}, in the spirit of the conformal reliability of
\cite{Li2026Conformal}.

% =====================================================================
\section{Discussion, Limitations, and Threats to Validity}
\label{sec:discussion}

\textbf{No empirical evidence yet.} This draft reports a method and a protocol. It
contains no measured results, and no accuracy, F1, or error figure anywhere in the text
should be read as one. The only numbers that appear are those quoted from prior work and
attributed to it.

\textbf{Novelty risk.} The retrieval-for-adaptation idea sits close to two literatures:
prototype/few-shot matching \cite{Yin2022FewSense,Zhao2024CrossFi,Yang2023FewCS} and
wireless RAG \cite{Nazar2024Enwar,Nazar2026Enwar2}. The distinguishing claim --- spatially
indexed memory conditioning a sensing head, updated label-free --- must survive a fresh
search and a carefully worded related-work section. If a concurrent paper makes the same
move, the fallback contribution is the disentanglement-plus-retrieval interaction,
which is testable independently.

\textbf{Disentanglement assumptions.} Perfect separation of content and domain is not
achievable in general; residual leakage is expected, and probing diagnostics are designed to
measure it rather than assert it. If $\zd$ leaks gesture information, the retrieval
conditioning becomes a soft shortcut that could inflate in-domain scores without
transferring.

\textbf{Retrieval realism.} The memory must be populated from captures that are cheap to
obtain, otherwise the approach relocates the data-collection burden rather than removing
it. The protocol must therefore report exactly what a new entry costs (duration of
unattended capture, number of devices, whether the space must be empty), and compare that
cost against the labelled data the baselines need.

\textbf{Privacy.} CSI is frequently preferred over cameras because it does not image
people \cite{Radwan2025Survey}. The geometric descriptors we store are a property of the
space, not of a person, and the memory must be specified so that it cannot be inverted into
person-identifying information. This is a design constraint, not an afterthought.

% =====================================================================
\section{Conclusion}
\label{sec:conclusion}

We propose \method{}, which treats the environment in cross-domain Wi-Fi CSI sensing as
retrievable context rather than nuisance: a self-supervised objective factorises a CSI
window into content and domain codes, and a spatially indexed memory conditions the
classifier on the space in which the measurement was taken, so that a new room is handled
by inserting an entry instead of retraining a model. We have stated the method, the
falsifiable claim--experiment mapping, and the baseline set that would establish or refute
it. The empirical work remains to be executed, and no result in this manuscript has been
measured. The immediate next step is to run the source-domain training and the memory-free
ablation on Widar3.0, because those two results decide whether the rest of the paper is
worth writing.

\bibliographystyle{IEEEtran}
\bibliography{references}

\end{document}
